\chapter{第一代总览：一个指纹，一个路标}
\label{ch:gen1-big}

\begin{genonebox}[第一代 FastCDC Stream 一句话]
从当前位置开始，\textbf{每读一个新字节}就更新一次指纹 $h$；  
当 $h$ 的\textbf{某些二进制位全是 0}（像中了「路标」），就在这儿切一刀。
\end{genonebox}

\section{用「高速公路出口」理解}

\begin{figure}[htbp]
\centering
\begin{tikzpicture}[scale=0.95]
  \draw[thick, Gen1Blue!60] (-5,0) -- (5,0);
  \node[font=\small] at (0,-0.6) {文件字节流（单向开过去）};
  \foreach \x in {-4,-2,0,2,4} {
    \draw[Gen1Blue, thick] (\x,0.1) -- (\x,0.5);
  }
  \node[gen1, minimum width=1.8cm] at (-4,1.2) {读字节};
  \node[gen1, minimum width=1.8cm] at (-2,1.2) {更新 $h$};
  \node[gen1, minimum width=1.8cm] at (0,1.2) {查路标};
  \node[gen1, minimum width=1.8cm] at (2,1.2) {命中？};
  \node[gen2, minimum width=1.8cm] at (4,1.2) {切！};
  \draw[arr] (-4,0.5) -- (-4,0.9);
  \draw[arr] (-2,0.5) -- (-2,0.9);
  \draw[arr] (0,0.5) -- (0,0.9);
  \draw[arr] (2,0.5) -- (2,0.9);
  \draw[arr, CutRed] (4,0.5) -- (4,0.9);
\end{tikzpicture}
\caption{循环四步：读 → 滚指纹 → 看路标 → 切或不切。}
\label{fig:gen1-loop}
\end{figure}

\section{第一代用到的四样东西}

\begin{table}[htbp]
\centering
\caption{第一代核心组件（先记名字，后面逐章揉碎）}
\begin{tabular}{@{}ll@{}}
\toprule
\textbf{名字} & \textbf{干什么} \\
\midrule
Gear 表 & 256 个「魔法数」，每个字节查一次 \\
滚动窗口 & 最近 $w$ 个字节（通常 64）决定指纹 \\
Mask（路标） & 一个位模式：$(h \ \& \ \text{mask}) == 0$ 就切 \\
min / max & 切太早、切太晚的护栏 \\
\bottomrule
\end{tabular}
\end{table}

\section{伪代码（整章只有这一处，别怕）}

\begin{codebox}[第一代切一块 chunk 的骨架]
\begin{lstlisting}
pos = chunk_start
scan_start = pos + min_size
cut_limit  = min(pos + max_size, file_end)

h = 预热(文件, scan_start, 窗口 w)   // 先「暖机」
for i from scan_start to cut_limit-1:
    if (h & mask) == 0:               // 路标命中
        cut_at = i
        break
    b = 文件[i]
    h = 滚动更新(h, b, 窗口 w)        // 下一章细讲
if 没找到:
    cut_at = cut_limit                // max 强制切
\end{lstlisting}
\end{codebox}

\begin{stepbox}[Step 3 — 第一代核心]
只有一个指纹 \keyterm{$h$}，只有一个路标 \keyterm{mask}。  
后面第二代只是把「一个变两个」，护栏和循环结构\textbf{不变}。

\end{stepbox}

\section{模拟示例：伪代码第一轮}

对统一玩具文件，第一轮循环状态：

\begin{simbox}[骨架代码代入数值]
\begin{lstlisting}
pos = 0
scan_start = 0 + 2 = 2
cut_limit  = min(0 + 8, 12) = 8
h = 暖机(...) = 18          // 见第 5 章
// i = 2:
if ((18 & 3) == 0)  // 18&3=2 != 0 -> 不切
h = 滚动(18, 文件[2]=0x30) = 20
// i = 3:
if ((20 & 3) == 0)  // 命中 -> cut_at = 3
\end{lstlisting}
\end{simbox}
